\begin{python0}
from solutions import *; clear()
\end{python0}
\sectionthree{Casting}

Recall that from CISS240, you can typecast an \texttt{int} to a \texttt{double}, a \texttt{double} to an \texttt{int}, a \texttt{char} to an \texttt{int}, an \texttt{int} to a \texttt{char}, etc.:
\begin{console}
int i = 42;
double d = 3.14;
char c = 'A';
int flag = 1;

double j = double(i); // or (double) i
int e = int(d); // or (int) d
int f = int(c); // or (int) c
bool flag2 = bool(flag);

// ...
\end{console}

Recall from the notes on pointers, you cannot assign pointer values to a
pointer variable if the pointer types are not the same:
\begin{console}
int i = 42;
int * p = &i; // OK
double * q = &i; // BAD!!!
\end{console}

But in fact you can \ldots if you perform an explicit typecast:
\begin{console}
int i = 42;
int * p = &i; // OK
double * q = (double *)i; // OK ... but probably a
                          // bad idea ...
\end{console}
In general, it's a good idea to think carefully before you typecast.

The above typecast has been around since C (the parent language of C++
and also the ancestor language of many programming languages) and these
are available in C++. C++ has also introduced it's own casting
mechanisms through four different casting operators:

\begin{itemize}
\item
  \verb!static_cast!
\item
  \verb!const_cast!
\item
  \verb!reinterpret_cast!
\item
  \verb!dynamic_cast!
\end{itemize}

